\subsection{2.7 Der Übergang zur Projektionstheorie}

Die Analyse des universellen Stabilitätsfunktionals markierte einen entscheidenden Wendepunkt in der Entwicklung der PST.

Bis zu diesem Zeitpunkt lautete die grundlegende Fragestellung:

\[
\boxed{
\text{Kann der OPP aus sich selbst heraus eine vierdimensionale Raumzeit erzeugen?}
}
\]

Die numerischen Ergebnisse machten jedoch deutlich, dass diese Fragestellung wahrscheinlich falsch gestellt war.

Nicht weil die PST scheiterte.

Sondern weil möglicherweise eine stillschweigende Annahme hinterfragt werden musste.

%%%%%%%%%%%%%%%%%%%%%%%%%%%%%%%%%%%%%%%%%%%%%%%%%%%%%%%%%%%%%%

\section*{Der entscheidende Perspektivwechsel}

Die ursprüngliche Vorstellung war:

\[
\text{OPP}
\longrightarrow
\text{Raumzeit}.
\]

Das bedeutet:

Die Raumzeit sollte bereits als Eigenschaft des OPP vorhanden sein.

Die Simulationen sprachen jedoch dagegen.

Der rohe OPP blieb hochdimensional.

Optimierungen funktionierten nur dann gut, wenn bereits implizit Informationen über den Zielrang eingebaut wurden.

Daraus entstand eine völlig neue Sichtweise.

Vielleicht gilt stattdessen

\[
\boxed{
\text{OPP}
\longrightarrow
\text{Projektion}
\longrightarrow
\text{Raumzeit}.
}
\]

Die Raumzeit wäre dann keine Eigenschaft des OPP,

sondern eine Eigenschaft der Art und Weise,

wie ein Beobachter Informationen aus dem OPP gewinnt.

%%%%%%%%%%%%%%%%%%%%%%%%%%%%%%%%%%%%%%%%%%%%%%%%%%%%%%%%%%%%%%

\section*{Die Analogie zur Wahrnehmung}

Zur Motivation dieser Idee wurde eine einfache physikalische Analogie herangezogen.

Der Mensch erlebt die Welt ebenfalls niemals vollständig.

Unsere Sinnesorgane besitzen starke Einschränkungen.

Beispiele:

\begin{itemize}

\item

Das Auge nimmt nur einen winzigen Bereich des elektromagnetischen Spektrums wahr.

\item

Das Ohr registriert nur einen kleinen Frequenzbereich.

\item

Viele Tiere besitzen völlig andere Wahrnehmungsräume.

\end{itemize}

Die physikalische Realität ist also wesentlich größer als unsere Beobachtung.

Formal:

\[
\boxed{
\text{Wahrnehmung}
\neq
\text{Realität}.
}
\]

Warum sollte dies beim OPP anders sein?

%%%%%%%%%%%%%%%%%%%%%%%%%%%%%%%%%%%%%%%%%%%%%%%%%%%%%%%%%%%%%%

\section*{Die Rolle des Beobachters}

Hier entstand erstmals eine zentrale Hypothese der späteren PST:

Der Beobachter ist kein äußerlicher Zusatz,

sondern Bestandteil der Projektion.

Der OPP enthält möglicherweise eine enorme Zahl an Freiheitsgraden.

Ein realer Beobachter besitzt jedoch nur eine endliche Informationskapazität.

Dadurch entsteht zwangsläufig eine Informationsreduktion.

Formal:

\[
\mathcal{P}
:
\mathrm{OPP}
\rightarrow
\mathcal{O},
\]

wobei

\[
\mathcal{P}
\]

einen Projektionsoperator bezeichnet

und

\[
\mathcal{O}
\]

den beobachtbaren Zustandsraum.

%%%%%%%%%%%%%%%%%%%%%%%%%%%%%%%%%%%%%%%%%%%%%%%%%%%%%%%%%%%%%%

\section*{Quantisierung als Projektionsartefakt}

Bereits früher war im OP diskutiert worden,

dass Quantisierung möglicherweise keine fundamentale Eigenschaft der Natur ist.

Vielmehr könnte sie aus einer endlichen Auflösung des Beobachters entstehen.

Dies fügte sich nun überraschend gut in die neue Projektionstheorie ein.

Falls jede Beobachtung eine Informationsreduktion darstellt,

dann folgt unmittelbar:

\[
\boxed{
\text{Diskrete Messwerte können aus kontinuierlichen Strukturen entstehen.}
}
\]

Der OPP selbst müsste dann gar nicht quantisiert sein.

Erst die Projektion erzeugt diskrete Beobachtungen.

%%%%%%%%%%%%%%%%%%%%%%%%%%%%%%%%%%%%%%%%%%%%%%%%%%%%%%%%%%%%%%

\section*{Raumzeit als Projektion}

Noch überraschender war die folgende Überlegung.

Nicht nur Quantisierung,

sondern möglicherweise auch Raum selbst

könnte durch Projektion entstehen.

Dann wäre

\[
\boxed{
\text{Raumzeit}
=
\text{Beobachtungsstruktur}.
}
\]

Der OPP bliebe weiterhin hochdimensional.

Die vierdimensionale Raumzeit wäre lediglich die effektiv beobachtbare Geometrie.

%%%%%%%%%%%%%%%%%%%%%%%%%%%%%%%%%%%%%%%%%%%%%%%%%%%%%%%%%%%%%%

\section*{Verbindung zur Speziellen Relativität}

Diese Idee harmoniert bemerkenswert gut mit der Speziellen Relativitätstheorie.

Bereits dort existiert

kein absolutes Jetzt.

Verschiedene Beobachter besitzen unterschiedliche Gleichzeitigkeitsebenen.

Die PST erweitert diesen Gedanken.

Nicht nur Zeit,

sondern möglicherweise auch

Raum

ist beobachterabhängig.

Der OPP selbst bleibt dagegen

absolut.

%%%%%%%%%%%%%%%%%%%%%%%%%%%%%%%%%%%%%%%%%%%%%%%%%%%%%%%%%%%%%%

\section*{Neue Forschungsstrategie}

Mit diesem Perspektivwechsel änderte sich auch die gesamte numerische Methodik.

Statt den OPP selbst weiter zu verändern,

wurden nun verschiedene mathematische Projektionen untersucht.

Die zentrale Forschungsfrage lautete:

\[
\boxed{
\text{Welche Projektionen erzeugen stabile physikalische Strukturen?}
}
\]

Nicht jede Projektion besitzt dabei physikalische Bedeutung.

Gesucht wurden Projektionsoperatoren,

die

\begin{itemize}

\item

robust,

\item

reproduzierbar,

\item

parameterstabil,

\item

und möglichst universell

\end{itemize}

funktionieren.

%%%%%%%%%%%%%%%%%%%%%%%%%%%%%%%%%%%%%%%%%%%%%%%%%%%%%%%%%%%%%%

\section*{Kandidaten für Projektionsoperatoren}

Aus der modernen Datenanalyse existieren zahlreiche Verfahren,

welche hochdimensionale Strukturen in niedrigdimensionale Darstellungen überführen.

Zu den untersuchten Verfahren gehörten insbesondere:

\begin{itemize}

\item

Principal Component Analysis (PCA)

\item

Kernel Principal Component Analysis (Kernel PCA)

\item

Diffusion Maps

\item

Random Projection

\item

verschiedene Kernel-Familien

\end{itemize}

Diese Verfahren wurden bewusst gewählt,

weil sie keine physikalischen Annahmen enthalten.

Falls eines davon dennoch reproduzierbar physikalisch sinnvolle Strukturen erzeugt,

wäre dies ein starkes Indiz für die PST.

%%%%%%%%%%%%%%%%%%%%%%%%%%%%%%%%%%%%%%%%%%%%%%%%%%%%%%%%%%%%%%

\section*{Das Ziel der Projektion}

Die Projektionen sollten insbesondere beantworten,

ob sich aus demselben hochdimensionalen OPP

stabile Werte für

\[
d_{\mathrm{eff}}
\approx
4
\]

ergeben.

Wichtig war dabei:

Die Zahl Vier durfte

nicht

bereits Bestandteil des Algorithmus sein.

Nur dann besitzt das Ergebnis physikalische Aussagekraft.

%%%%%%%%%%%%%%%%%%%%%%%%%%%%%%%%%%%%%%%%%%%%%%%%%%%%%%%%%%%%%%

\section*{Der wissenschaftliche Gewinn}

Die Einführung der Projektionstheorie stellte den vielleicht größten konzeptionellen Fortschritt der gesamten PST dar.

Die Theorie löste sich damit endgültig von der Vorstellung,

Raumzeit müsse bereits fundamental existieren.

Stattdessen entstand folgendes Bild:

\[
\boxed{
\text{Existenz}
\rightarrow
\text{OPP}
\rightarrow
\text{Projektion}
\rightarrow
\text{Raumzeit}
\rightarrow
\text{Physik}.
}
\]

Raum und Zeit erscheinen damit nicht mehr als Grundbausteine der Wirklichkeit,

sondern als emergente Größen,

die erst im Akt der Projektion entstehen.

%%%%%%%%%%%%%%%%%%%%%%%%%%%%%%%%%%%%%%%%%%%%%%%%%%%%%%%%%%%%%%

\section*{Ausblick}

Die folgenden Kapitel dokumentieren die umfangreichen numerischen Untersuchungen dieser Projektionsoperatoren.

Dabei zeigte sich überraschenderweise,

dass nicht alle Verfahren gleichwertig sind.

Während einige Projektionen keinerlei stabile Strukturen erzeugten,

kristallisierten sich zwei Verfahren als besonders vielversprechend heraus:

\begin{enumerate}

\item Diffusion Maps

\item Kernel Principal Component Analysis (Kernel PCA)

\end{enumerate}

Diese beiden Methoden bilden den Ausgangspunkt der weiteren numerischen Entwicklung der PST und werden in den folgenden Kapiteln ausführlich untersucht.

